\documentclass[12pt,a4paper]{article}
\usepackage[utf8]{inputenc}
\usepackage[T1]{fontenc}
\usepackage{amsmath}
\usepackage{booktabs}
\usepackage{geometry}
\usepackage{hyperref}
\usepackage{cite}
\geometry{a4paper,top=2.5cm,bottom=2.5cm,left=2.5cm,right=2.5cm}
\hypersetup{colorlinks=true,linkcolor=blue,citecolor=blue,urlcolor=blue}

\title{\textbf{Water as Simultaneous Fuel, Moderator and Reaction Medium\\
for Proton--Electron Capture Chain Reactions\\
with Lithium-6 Energy Extraction}}
\author{Eymen Aydoğan}
\date{May 19, 2026}

\begin{document}
\maketitle

\begin{abstract}
We propose that ordinary water (H$_2$O) is an ideal reaction medium
for the PECCNAL proton--electron capture mechanism, serving
simultaneously as fuel source (hydrogen atoms), neutron moderator,
and reaction containment. Each water molecule provides 10 electrons
as potential capture targets (2 from hydrogen, 8 from oxygen),
offering electron density 11,000 times higher than hydrogen gas at
standard conditions. Produced neutrons are naturally moderated by
hydrogen nuclei in water before reaching a dissolved Lithium-6
target, which releases 4.78 MeV per neutron capture. The system
requires no cryogenic cooling, operates at room temperature, and
uses the most abundant liquid on Earth as its primary fuel.
\end{abstract}

\noindent\textbf{Keywords:} water medium, neutron moderation,
proton-electron capture, lithium-6, room temperature, LENR

\hrule\vspace{1em}

\section{Introduction}
Previous papers in this series~\cite{aydogan2026a} employed
hydrogen or deuterium gas as the reaction medium. Here we propose
that liquid water offers substantial advantages as a reaction medium
for the PECCNAL mechanism, combining high electron density, natural
neutron moderation, and near-zero material cost.

\section{Water as Reaction Medium}

\subsection{Electron Density Advantage}

\begin{table}[h!]
\centering
\caption{Electron density comparison across reaction media}
\vspace{0.5em}
\begin{tabular}{lcccc}
\toprule
\textbf{Medium} & \textbf{Density} & \textbf{e$^-$/molecule} & \textbf{Rel. density} \\
\midrule
H$_2$ gas (STP) & 0.09 kg/m$^3$ & 2 & 1$\times$ \\
H$_2$ gas (100 atm) & 9 kg/m$^3$ & 2 & 100$\times$ \\
Liquid H$_2$ & 70 kg/m$^3$ & 2 & 778$\times$ \\
\textbf{Liquid H$_2$O} & \textbf{1000 kg/m$^3$} & \textbf{10} & \textbf{111,000$\times$} \\
\bottomrule
\end{tabular}
\end{table}

The electron density advantage of water means a free proton
travelling through water encounters an electron target
111,000 times more frequently than in hydrogen gas at STP.

\subsection{Natural Neutron Moderation}

Produced neutrons must reach thermal energies for efficient
Li-6 capture. In water:
\begin{equation}
n + {}^1\text{H} \rightarrow {}^2\text{H} + \gamma
\quad \sigma_\text{thermal} = 0.33\,\text{barn}
\end{equation}

Water is the most effective neutron moderator known, used in
pressurized water reactors (PWR) worldwide. Neutrons thermalize
within centimetres of their production point, maximizing
Li-6 capture probability.

\subsection{Lithium-6 Dissolution}

Lithium carbonate (Li$_2$CO$_3$) dissolves in water, distributing
Li-6 uniformly throughout the reaction volume:
\begin{equation}
n + {}^6\text{Li} \rightarrow {}^4\text{He} + {}^3\text{H}
+ 4.78\,\text{MeV}
\end{equation}

\section{Reaction Sequence in Water Medium}

\begin{align}
p + e^-_\text{(H$_2$O)} &\rightarrow n + \nu_e
\quad -1.293\,\text{MeV} \tag{W1}\\
n + {}^1\text{H}_\text{(H$_2$O)} &\rightarrow {}^2\text{H} + \gamma
\quad +2.224\,\text{MeV} \tag{W2}\\
n + {}^6\text{Li}_\text{(dissolved)} &\rightarrow {}^4\text{He} + {}^3\text{H}
\quad +4.780\,\text{MeV} \tag{W3}
\end{align}

Net energy per cycle:
\begin{equation}
E_\text{net} = 2.224 + 4.780 - 1.293 = \boxed{+5.711\,\text{MeV}}
\end{equation}

\section{Practical Advantages}

\begin{itemize}
\item \textbf{No cryogenics}: operates at room temperature
\item \textbf{Self-shielding}: water absorbs most radiation
\item \textbf{Abundant fuel}: seawater contains $\sim$35 ppm Li
\item \textbf{Low cost}: water and Li$_2$CO$_3$ are inexpensive
\item \textbf{Scalable}: volume scales cubically with linear dimension
\end{itemize}

\section{Conclusion}
Water is proposed as the optimal reaction medium for PECCNAL,
offering 111,000$\times$ higher electron density than hydrogen gas,
built-in neutron moderation, and compatibility with dissolved Li-6.
Net energy output is 5.711 MeV per cycle with no cryogenic
requirements.

\section*{Acknowledgments}
The author thanks the nuclear and LENR physics communities.

\begin{thebibliography}{9}
\bibitem{aydogan2026a} E. Aydoğan, ``Coulomb-Barrier-Free Neutron Production,'' arXiv preprint, 2026.
\bibitem{pdg2022} R. L. Workman et al. (PDG), \textit{Prog. Theor. Exp. Phys.}, 083C01, 2022.
\end{thebibliography}
\end{document}
